\section{Preliminaries}
\label{sec:preliminaries}

\subsection{Optimal Transport Conditional Flow Matching}

Flow matching~\cite{lipman2023flow} learns a time-dependent velocity field $v_\theta: \RR^d \times [0,1] \to \RR^d$ that generates a probability path $p_t$ interpolating between a prior distribution $p_0$ (e.g., standard Gaussian) and the data distribution $p_1$. A sample $x_0 \sim p_0$ is transported to $x_1 \sim p_1$ by integrating the ODE:
\begin{equation}
    \frac{\mathrm{d}x_t}{\mathrm{d}t} = v_\theta(x_t, t), \quad x_0 \sim p_0, \quad t \in [0, 1].
    \label{eq:flow_ode}
\end{equation}

The Conditional Flow Matching (CFM) objective trains $v_\theta$ by regressing against the conditional velocity field $u_t(x | x_1)$ that generates the conditional probability path $p_t(x | x_1)$:
\begin{equation}
    \LL_{\text{CFM}} = \EE_{t \sim \mathcal{U}[0,1],\, x_1 \sim p_1,\, x_t \sim p_t(\cdot|x_1)} \left\| v_\theta(x_t, t) - u_t(x_t | x_1) \right\|^2.
    \label{eq:cfm_loss}
\end{equation}

With Optimal Transport CFM (OT-CFM)~\cite{tong2024improving}, the conditional path is a linear interpolation with OT coupling:
\begin{equation}
    x_t = (1 - t)\, x_0 + t\, x_1, \quad u_t(x_t | x_1) = x_1 - x_0,
    \label{eq:ot_interpolation}
\end{equation}
where $(x_0, x_1)$ are matched via the OT plan $\pi^* = \argmin_\pi \EE_{(x_0,x_1)\sim\pi}\|x_0 - x_1\|^2$. This yields straighter flow trajectories and faster convergence than standard CFM.

At inference, actions are generated by solving the ODE via Euler integration with $K$ steps:
\begin{equation}
    x_{t+\Delta t} = x_t + v_\theta(x_t, t) \cdot \Delta t, \quad \Delta t = 1/K.
    \label{eq:euler_sampling}
\end{equation}

\subsection{VITA Architecture}
\label{sec:vita_arch}

VITA~\cite{gao2026vita} instantiates flow matching for visuomotor policies with a latent-space formulation. The architecture consists of four components (Figure~\ref{fig:architecture}, gray region):

\paragraph{Observation Encoder $\FF_\phi$.} A ResNet-18~\cite{he2016deep} backbone extracts visual features from $N$ camera images, which are pooled and concatenated with the proprioceptive state $s_t \in \RR^{d_s}$:
\begin{equation}
    \bm{o}_t = \FF_\phi\!\left(\{I_t^{(i)}\}_{i=1}^N,\; s_t\right) \in \RR^{d_o}, \quad d_o = 512N + d_s.
    \label{eq:obs_encoder}
\end{equation}
A linear projection maps this to the latent space: $\bm{z}_t^{\text{obs}} = W_{\text{enc}}\, \bm{o}_t \in \RR^{d_z}$.

\paragraph{Action Autoencoder.} A 1D-CNN encoder $E_\psi$ maps action chunks $\bm{a}_{t:t+H} \in \RR^{H \times d_a}$ to latent vectors $\bm{z}^{\text{act}} = E_\psi(\bm{a}_{t:t+H}) \in \RR^{d_z}$, and an MLP decoder $D_\psi$ reconstructs actions: $\hat{\bm{a}}_{t:t+H} = D_\psi(\bm{z}^{\text{act}})$.

\paragraph{Flow Network $G_\theta$.} An MLP with FiLM~\cite{perez2018film} time conditioning predicts the velocity field. Each layer $l$ applies:
\begin{equation}
    \bm{h}^{(l)} = \bm{h}^{(l-1)} + \gamma_t^{(l)} \odot \text{MLP}\!\left((\bm{1} + \bm{\alpha}_t^{(l)}) \odot \text{LN}(\bm{h}^{(l-1)}) + \bm{\beta}_t^{(l)}\right),
    \label{eq:film_layer}
\end{equation}
where $(\gamma_t, \bm{\alpha}_t, \bm{\beta}_t) = f^{(l)}(\text{SinEmb}(t))$ are FiLM parameters derived from a sinusoidal embedding of the flow timestep $t$.

\paragraph{Training Objective.} VITA uses a composite loss:
\begin{equation}
    \LL_{\text{VITA}} = \LL_{\text{flow}} + \lambda_{\text{enc}} \LL_{\text{enc}} + \lambda_{\text{fld}} \LL_{\text{fld}} + \lambda_{\text{con}} \LL_{\text{con}},
    \label{eq:vita_loss}
\end{equation}
where $\LL_{\text{flow}}$ is the OT-CFM loss (Eq.~\ref{eq:cfm_loss}), $\LL_{\text{enc}}$ is the action autoencoder reconstruction loss, $\LL_{\text{fld}}$ is the flow-decoded reconstruction loss, and $\LL_{\text{con}} = \|\bm{z}^{\text{act}}_{\text{pred}} - \bm{z}^{\text{act}}\|^2$ is a latent consistency loss. Crucially, the flow starts from the observation latent $\bm{z}_t^{\text{obs}}$ rather than random noise, establishing a direct vision-to-action mapping.

\subsection{FiLM Conditioning}
\label{sec:film}

Feature-wise Linear Modulation (FiLM)~\cite{perez2018film} conditions a neural network layer on an external signal $c$ through learned affine transformations:
\begin{equation}
    \text{FiLM}(\bm{h}; c) = (\bm{1} + \bm{\alpha}_c) \odot \bm{h} + \bm{\beta}_c,
    \label{eq:film}
\end{equation}
where $\bm{\alpha}_c, \bm{\beta}_c = f(c)$ are predicted scale and shift parameters. FiLM is widely used in conditional generative models for injecting timestep, class, or style information. In standard flow networks, FiLM conditions only on the flow timestep $t$. Our Dual FiLM extends this to jointly condition on flow timestep and task progress through separate modulation pathways.
